\documentclass[10pt,letterpaper,twocolumn]{article}
% EMC2 (AAAI-27 workshop) format emulation of the AAAI two-column author kit
% (official aaai2027.sty not installed on this machine; geometry matches the kit:
% letter paper, 7.0in text width, 0.375in column separation, Times 10pt).
\usepackage[letterpaper,textwidth=7.0in,textheight=9.0in,top=0.75in,left=0.75in,columnsep=0.375in]{geometry}
\setlength{\textfloatsep}{10pt plus 2pt minus 2pt}
\setlength{\floatsep}{8pt plus 2pt minus 2pt}
\setlength{\intextsep}{8pt plus 2pt minus 2pt}
\usepackage{times}
\usepackage{booktabs}
\usepackage{amsmath}
\usepackage[hidelinks]{hyperref}
\usepackage{natbib}
\emergencystretch=2.5em
\title{When Do Parameter Budgets Misprice CPU Inference?\\ A Pre-registered Limit Study of Byte-Traffic Pruning Allocation}
\author{Prabhu Rangaraju\\ Independent Researcher}
\date{}
\begin{document}
\maketitle


\begin{abstract}
Pruning budgets are almost always stated in
parameter count, yet CPU decode cost is plausibly driven by the bytes
moved per token: weight bytes streamed each step plus KV-cache bytes
read at the occupied context. We test, under a pre-registered
protocol, whether allocating a $\sim$25\% structured prune of
Llama-3.2-3B by an analytical per-token byte-traffic model beats
parameter-count allocation, at fixed CPU thread caps
(2 and 4 threads, both pinned to 2 cores). Eleven
allocations search one shared unit space (layers, FFN channels, KV
groups) and differ only in the cost model choosing the cut.
The tightened evidence does not support the negative
result the protocol was built to deliver. A teacher-forced validity check of
every exported variant (Amendment A3) passes 1 of 11:
only V-depth-lite's GGUF stays within the frozen export-fidelity
tolerances, so the pre-registered 10-variant allocation contest
(six stage-1-valid plus four amendment variants)
cannot be fitted on a D4-valid primary set, and the registered
head-to-head has both arms invalid. The frozen local contest over
the original six stage-1-valid variants, reported with that caveat:
the traffic model does not beat the parameter model by the registered margin at either cap ($\Delta$AIC $+2.21$ at 2 threads, $-19.11$ at 4), and
the predeclared byte allocation does not reliably meet the
registered standard (paired ratios 0.917 / 0.986 at 2 threads,
0.818 / 0.817 at 4), though in a post-hoc sensitivity excluding the saturated
16384 context -- an exclusion requested in review, not
pre-registered -- it does win the 4-thread contest. Quality is a
context axis only: the frozen quality scorer is void (it scores
raw-completion continuations, not answers), and the Amendment-4
first-answer re-score of the frozen outputs gives BASE 85.00
(34/40) and the sole D4-valid variant 55.00 (22/40). The contribution is
therefore a boundary, drawn twice over: 10 of 11 allocations
fail the export-fidelity standard under the joint
prune-and-quantise pipeline (the two steps are applied jointly
and never separated), and
the traffic model's own predicted speedups run systematically
ahead of what was measured. Traffic-aware pruning as
such sits outside it --- pruning by \emph{measured} latency is
untouched by anything we report. Scope: every timing comes from
a shared 2-core virtual machine (tenant-load variance is a
formal limitation); the outcome is CPU time, not energy.
\end{abstract}

\section{Introduction}
Structured pruning removes whole layers, channels, or attention
components so that a smaller model runs faster on real hardware.
The budget for such a prune is conventionally a parameter count:
remove $x$\% of the weights. Parameter count is, however, a proxy.
At batch size 1, autoregressive decode is widely treated as
memory-bound: each step must stream the model's weights once and
read the accumulated key--value (KV) cache, whose size grows with
context. A cost model based on per-token byte traffic --- weight
bytes streamed per step plus KV bytes read at the occupied context
--- is therefore an attractive alternative: it prices KV-heavy
allocations differently from parameter-heavy ones, and it predicts
that two allocations with identical parameter counts can have
different decode costs, especially at long contexts.

This intuition has been used to \emph{analyse} inference
\citep{rooflinebench2026} and named as an untested extension by
pruning work that itself budgets in parameters
\citep{grasprune2026}, but we find no prior controlled test of the
allocation question (literature search, October 2026):
\emph{within one matched-budget pruning
family, does choosing the cut by byte traffic produce a faster
model, and does the traffic model predict held-out allocations'
speed better than parameter count?}

We answer with a pre-registered, single-model, single-runtime
experiment in which the two cost models are the only things that
differ between arms. Export-validity failures restrict the
results to a subset of the design space; this paper reports
that narrower result honestly:
\begin{enumerate}
\item \textbf{Validity.} A teacher-forced check of every variant's
exported model passes 1 of 11 (Section~\ref{sec:validity}). The
pre-registered allocation contest cannot be fitted on a D4-valid
primary set; the registered head-to-head has both arms invalid.
\item \textbf{Frozen contest, with caveat.} On the original six
stage-1-valid variants, the traffic model does not out-predict the
parameter model (Section~\ref{sec:contest}), and the predeclared
byte allocation fails the pre-registered $\geq$10\% paired-advantage
criterion at both thread caps (Section~\ref{sec:hurdle}). Five of
those six fail the amended (D4) validity check on rebuilt models, so the contest is frozen evidence with that caveat.
\item \textbf{Model calibration.} Where the traffic model predicts
23--33\% advantages at cap 2, observation delivers 1--8\%; the
decision rule shares the blame, but the model itself is
faulty.
\item \textbf{Quality as context only.} The quality evaluation ran; its frozen scorer is void, and the Amendment-4 first-answer re-score is reported as a context axis, never a verdict criterion (Section~\ref{sec:quality}).
\end{enumerate}
We also report selection regret descriptively
(Section~\ref{sec:regret}): the byte-greedy pick was in raw terms a
decent pick, which makes the model's inaccuracy more informative.

\section{Related work}
\textbf{Measured-latency pruning.} ZipLM prunes by target-runtime
latency lookup tables rather than by size \citep{ziplm2023}; HAPE ranks structures by measured on-device
latency differences for Llama-2-7B on a Xeon CPU, the closest
deployment precedent to ours \citep{hape2025}. Both sidestep
analytical cost models altogether, and neither asks whether an
analytical traffic model adds predictive value over parameters
on unseen allocations.
\textbf{Structured pruning methods.} LLM-Pruner removes coupled
structures by gradient-estimated importance and retrains
\citep{llmpruner2023}; SliceGPT cuts dimensions in a rotated
space \citep{slicegpt2024}; Wanda prunes by weight magnitude
times activation norm \citep{wanda2024}. Our variants span the
same axes (layers, FFN width, KV heads) at a matched $\sim$25\%
removal, but this paper proposes no pruning method: the treatment
is the cost model that chooses the cut.
\textbf{Joint FFN--KV allocation.} GRASPrune allocates a single
budget jointly over FFN channels and KV groups, but states that
``we use parameter count as the budget proxy''; a bandwidth-aware
cost is proposed in its \S3.1 and named as a limitation, not
evaluated \citep{grasprune2026}. Our experiment is, to our
knowledge, the controlled test of that extension.
\textbf{KV-targeted pruning.} KVPruner prunes KV channels and
reports that equal parameter counts can yield different speeds,
but its inference-overhead evaluation disables caching, so its
numbers do not speak to cached decode \citep{kvpruner2024}.
\textbf{Traffic and roofline models.} RooflineBench defines
per-token traffic as weight bytes plus active KV reads/writes ---
the formula family we use --- and benchmarks it across on-device
hardware including Llama-3.2-1B, without a pruning-allocation test
\citep{rooflinebench2026}. Analytical models separate weight and
KV memory costs \citep{pope2023}; CPU characterisation and
kernel-level simulation both find that a single bandwidth
constant cannot capture decode performance
\citep{iiswc2024,kernelsight2026}.
\textbf{This venue.} EMC$^2$'s own program prices the same
quantities: KV-cache retention as a carbon tradeoff
\citep{bassi2026}; NF4 dequantisation as a first-order cost
\citep{qi2026}; efficiency by composite index \citep{parihar2026};
operation fusion, which changes performed traffic, as a speedup
\citep{salmani2025}. Our design asks whether the simplest
analytical traffic model earns its keep for one decision: where
to cut. The topic list names ``limit studies'' and ``bottleneck
analysis'' explicitly; this paper is one.

\section{Method}
\subsection{Variants and cost models}
We prune Llama-3.2-3B-Instruct (GGUF Q4\_K\_M; KV cache fp16) by
$\sim$25\% of parameters. Eleven allocations search one shared
search space (whole layers, FFN channels, KV groups); allocations
differ only in the cost model that prices each unit: parameter
count, marginal byte traffic $\Delta B$ at the design contexts, or
a mix of the two. Final parameter counts span 2.398--2.442B
(1.84\% spread); tensor-payload bytes for the original seven span
1.526--1.552 GB; parameter count enters all analyses as a
covariate, so exact equality is not load-bearing.
SIMD-alignment constraints bind all variants (FFN multiples of
256, hidden widths multiples of 64, head dimension 128). The six
originally stage-1-valid variants are the frozen contest set; the four amendment
variants (A1) extend it to ten for the amended contest.

Both candidates are produced by applying their cost model to the
\emph{identical} candidate space, so their differing composition
is the treatment, not a confound: the parameter model spends the
whole budget on FFN width (layers and KV heads are expensive in
parameters per unit), while the byte model spends it on layers
and KV heads (expensive in traffic per unit). The remaining
variants span intermediate selections of the same space, so the
prediction contest (Section~\ref{sec:contest}) is over
allocations, not only over the two endpoints.

\begin{table}[ht]
\centering\scriptsize\setlength{\tabcolsep}{4pt}
\begin{tabular}{@{}llllll@{}}
\toprule
Variant & Layers & KV-heads & Params (B) & Allocation rule & D4 \\
\midrule
V-width & 28 & 224 & 2.442 & params (FFN trim) & FAIL \\
V-depth & 21 & 168 & 2.409 & 7 layers dropped & FAIL \\
V-depth-lite & 24 & 192 & 2.414 & 4 layers dropped & PASS \\
V-depth-alt & 21 & 168 & 2.409 & 7 layers, alt.\ positions & FAIL \\
V-KV & 28 & 112 & 2.398 & KV heads 8$\rightarrow$4 & FAIL \\
V-KV-lite & 28 & 168 & 2.442 & KV heads 8$\rightarrow$6 & FAIL$^*$ \\
V-mixed & 23 & 92 & 2.420 & byte-greedy mix & FAIL \\
V-shallow & 26 & 208 & 2.398 & 2 top layers dropped & FAIL \\
V-kv2 & 28 & 56 & 2.420 & KV heads 8$\rightarrow$2 & FAIL \\
V-mid & 23 & 138 & 2.402 & 5 mid layers, KV 8$\rightarrow$6 & FAIL \\
V-front & 24 & 192 & 2.414 & 4 front layers dropped & FAIL \\
\bottomrule
\end{tabular}
\caption{The eleven allocations. V-width is the parameter-count
candidate and V-mixed, predeclared before measurement, is the
byte-traffic candidate. Parameter counts verify the matched budget
(2.398--2.442B, 1.84\% spread). D4 is the amended validity verdict
(Section~\ref{sec:validity}). $^*$V-KV-lite: the original build was
export-invalid (top-1 0.000); the amended rebuild fails D4 on KL by 0.005.}
\end{table}

The traffic model prices a decode step at occupied context $c$ as
$B(c) = B_{\text{weights}} + B_{\text{KV}}(c)$, with
$B_{\text{KV}}(c) = 2(\sum \text{KV-heads})\cdot128\cdot2\text{B}\cdot c$. For the unpruned base this gives KV
shares of total traffic of $\approx$24\% at 4k, 39\% at 8k, and
56\% at 16k context --- the long-context regime is where a traffic
model has room to differ from parameter count, so contexts below 4096 are reported but the registered contest uses $C \geq 4096$.

\subsection{Protocol}
The protocol, criteria, and analysis definitions were
pre-registered before any stage-1 measurement (study record
SPEC v0.3, 2026-10-23; runtime amendments were logged in the
same record and are summarised where they bear on the results).
Timing uses llama.cpp (\texttt{llama-server}) on a 2-core virtual
machine (2 vCPUs of an AMD EPYC 9D25 host). Thread caps are llama.cpp thread counts ($-t$ 2 and
$-t$ 4) under a \texttt{taskset -c 0,1} pinning to two cores; ``cap 4''
is four threads on two cores --- an oversubscribed cap, not a
four-core result --- labelled as such throughout. Each measurement
cell (variant $\times$ cap $\times$ context $\in$
\{512, 2048, 4096, 6144, 8192, 16384\}) runs in a fresh server
process (``virgin'' protocol, adopted after a slot-history
contamination artifact was caught and the 30 pre-protocol reps
excluded), with prefill-only primes, 5 warmup and 5 recorded reps of 128 decode steps on fixed token IDs.
All decoding is single-sequence (batch size 1); the primary metric is CPU seconds per token, with per-rep host load and context switches recorded. In total 521 virgin
reps are banked, of which 451 belong to the six primary variants
(see Section~\ref{sec:validity}); the contest fits use the
$C\geq4096$ subsets (135 reps at cap 2, 136 at cap 4).

\subsection{Pre-registered criteria}
Three frozen tests decide the claim, with one descriptive analysis
alongside them. (i) \emph{Prediction contest:}
regressions of log CPU-s/token on either parameter count or
modeled byte traffic (each with a layer-count term), judged on
leave-one-variant-out (LOVO) relative error over $C\geq4096$ reps;
the traffic model wins only if it has lower LOVO error
\emph{and} $\Delta$AIC (parameters $-$ traffic) $> 10$. LOVO
relative error is the mean, over the held-out variant's reps, of
$|\text{pred}-\text{obs}|/\text{obs}$; every held-out variant is an extrapolation to an unseen architecture, so absolute values are large for both models and the contest is comparative by design.
(ii) \emph{Registered hurdle:} in genuine temporally-paired blocks
(five blocks per cap $\times$ \{8192, 16384\}, adjacent
fresh-process cells, randomised order), the upper 95\% confidence
bound of the paired log-time ratio V-mixed/V-width must be
$< 0.90$ at both contexts per cap. (iii) \emph{Validity (A3):}
every variant's exported GGUF is checked teacher-forced against
its own HF pruned twin (144 positions); a variant failing the
frozen margins is excluded from primary fits, with a fit including the failed variants as sensitivity (Section~\ref{sec:validity}).
(iv) \emph{Selection regret} (post-hoc and non-registered;
descriptive only, computed at manuscript stage from banked data):
for each cell, regret is a selector's mean time minus the
fastest primary mean, relative to that mean.
The amended campaign (Amendments 1--3, 2026-11-04) adds four
variants, a quality evaluation, the all-variant validity check,
equivalence/power analyses, a FLOP-count comparator arm, and a
multicore timing arm. The amended timing (Kaggle, a public
cloud-compute platform; campaign K2) did not complete: the run
was cancelled at Kaggle's 12-hour session cap mid-way through
the fourth variant, and the harvested partial measurements
cover 46 of 120 cells over 4 of 10 variants (no verdict file
was produced). The analyses consuming the amended timing (the
including-fit sensitivity, TOST, power, the no-16384 secondary,
and the FLOP arm, Amendment 2 item B1) therefore remain unrun
and are the named remaining shortfall, outside this
submission's scope; no substitute data are used for them.
Amendment 4 (2026-11-05) corrects the quality scorer; see
Section~\ref{sec:quality}.

\section{Results}
\subsection{Validity: 1 of 11 exports survives}
\label{sec:validity}
\begin{table}[ht]
\centering\scriptsize\setlength{\tabcolsep}{4pt}
\begin{tabular}{@{}llllll@{}}
\toprule
Variant & HF NLL & dNLL & KL & top-1 & D4 \\
\midrule
BASE & 2.968 & $-$0.009 & 0.200 & 0.951 & ref \\
V-width & 8.498 & $+$0.003 & 0.446 & 0.660 & FAIL \\
V-depth & 4.376 & $+$0.074 & 0.275 & 0.868 & FAIL \\
V-depth-lite & 3.818 & $+$0.005 & 0.263 & 0.889 & PASS \\
V-depth-alt & 4.225 & $-$0.011 & 0.468 & 0.826 & FAIL \\
V-KV & 7.534 & $-$0.032 & 0.937 & 0.681 & FAIL \\
V-KV-lite & 3.614 & $+$0.012 & 0.404 & 0.903 & FAIL \\
V-mixed & 8.488 & $+$0.102 & 0.860 & 0.667 & FAIL \\
V-shallow & 4.576 & $-$0.064 & 0.389 & 0.833 & FAIL \\
V-kv2 & 7.811 & $-$0.165 & 1.043 & 0.743 & FAIL \\
V-mid & 4.207 & $+$0.024 & 0.590 & 0.833 & FAIL \\
V-front & 4.752 & $+$0.008 & 0.638 & 0.826 & FAIL \\
\bottomrule
\end{tabular}
\caption{Teacher-forced validity (A3; 144 positions; Kaggle
campaign K1). dNLL, KL and top-1 compare each variant's GGUF
against its own HF pruned twin, so D4 tests the export, not the
pruning: a variant fails D4 when its GGUF diverges from its HF
twin (quantisation divergence), whatever its pruning damage (the
HF NLL column). (failing variants' KL is 1.4--5.2$\times$ BASE's 0.200.)
Limits: dNLL $\leq$ 0.041 (one-sided: dNLL(V) $\leq$
dNLL(BASE) $+$ 0.05), KL $\leq$ 0.399, top-1 $\geq$ 0.851.
Two failures are marginal (V-KV-lite KL by 0.005; V-shallow
top-1 by 3 positions). All 11 verdicts were recomputed
from the harvested K1 files and match; the K1 raw-output pull
has since landed and satisfies the quality-output provenance
(Section~\ref{sec:quality}).}
\end{table}
The check runs on Kaggle (campaign K1; pinned base snapshot and
llama.cpp 836d571; 7,956 s); the seven original variant GGUFs
are byte-identical (sha256) to the test-9 clean-room rebuilds,
whose packed bytes match the frozen analysis exactly. The K1 BASE GGUF
differs from the frozen BASE by 128 bytes (presumed metadata;
a tensor-byte comparison is not possible from the raw pull,
which contains no .gguf files), so
transfer of the K1 verdicts to the frozen contest's artifacts
is conditional on that comparison; the verdicts are internal
to K1's own BASE reference and unaffected. The dNLL bound is
one-sided by the frozen rule (a GGUF scoring better than its HF
twin passes), so V-shallow ($-$0.064) and V-kv2 ($-$0.165)
fail D4 on top-1 and KL, not dNLL. The 1-of-11 count is
margin-sensitive: KL $\leq$ 0.45 admits V-KV-lite (2 of 11);
relaxing all three limits ($|$dNLL$|$ $\leq$ 0.07, KL $\leq$ 0.45,
top-1 $\geq$ 0.83) admits V-shallow too (3 of 11); that most
variants fail is robust to these relaxations. Two readings matter. First, the D4 rule as frozen is
binary and A3 applies it to primary-fit membership: the primary
fit set is \{V-depth-lite\} alone, so the 10-variant contest
cannot be fitted as a primary analysis, and a fit including all
variants is the sensitivity analysis. Second, the pattern is itself the
result: pruning damage (HF NLL) and quantisation divergence move
together --- the four variants with the worst pruning damage
(V-width, V-mixed, V-kv2, V-KV; NLL $+$4.6 to $+$5.5) are also the four with KL 2.2--5.2$\times$ BASE's. The registered hurdle's two arms both fail
D4; its timings are retained per the frozen rule, reported as timings of artifacts diverging from their intended pruned models beyond tolerance.
V-KV-lite deserves a correction: its stage-1 invalidity (top-1
0.000, a converter bug) described the original build; the amended
rebuild gives top-1 0.903 and dNLL $+$0.012 --- valid in all but a marginal KL failure (0.404 against 0.399), so the amended verdict is a marginal failure, not an invalid export; the original exclusion was correct at the time.

\subsection{Registered hurdle: not met (both arms invalid)}
\label{sec:hurdle}
\begin{table}[ht]
\centering\small
\begin{tabular}{@{}lllll@{}}
\toprule
Cap & Ctx & Paired ratio & 95\% CI & Traffic-predicted \\
\midrule
2 & 8192 & 0.917 & [0.824, 1.020] & 0.771 \\
2 & 16384 & 0.986 & [0.934, 1.041] & 0.673 \\
4 & 8192 & 0.818 & [0.795, 0.842] & 0.771 \\
4 & 16384 & 0.817 & [0.609, 1.096] & 0.673 \\
\bottomrule
\end{tabular}
\caption{Paired V-mixed/V-width time ratios (five retained
pairs per cell, from the paired-block manifest). The
traffic model's predicted ratio is cap-independent (it prices
bytes, not threads). The
hurdle (upper bound $<0.90$ at both contexts) is met in neither cap:
cap 2 fails: the upper bounds exceed 0.90 and the intervals include 1; cap 4 passes at 8192 but is too imprecise at 16384. Caveat: both arms fail D4
(Section~\ref{sec:validity}); these are timings of the exported
artifacts, not of validated pruned models.}
\end{table}
At 4 threads the byte candidate wins ($\approx$0.82); at 2 threads it does not (0.92/0.99, CIs include 1). Nowhere does it meet the registered standard --- and both arms fail D4, so these paired ratios compare exported artifacts that diverge from their intended pruned models beyond tolerance, not validated pruned models. \subsection{Prediction contest: traffic does not win (frozen, caveated)}
\label{sec:contest}
\begin{table}[ht]
\centering\footnotesize\setlength{\tabcolsep}{4pt}
\begin{tabular}{@{}llll@{}}
\toprule
Cap & Params err & Traffic err & $\Delta$AIC (P$-$B) \\
\midrule
2 & 0.743 & 0.689 & $+2.21$ (needed $>10$) \\
4 & 0.284 & 0.538 & $-19.11$ (params win) \\
\bottomrule
\end{tabular}
\caption{Held-out LOVO prediction error, $C\geq4096$ (135 reps at
cap 2, 136 at cap 4), frozen six-variant primary fit. Including
V-KV-lite (invalid in stage 1) as a sensitivity gives the same outcome
($\Delta$AIC $+3.37$ at cap 2, $-19.56$ at cap 4). Caveat: five of these six
stage-1-valid variants fail the amended (D4) validity check on rebuilt GGUFs
(Section~\ref{sec:validity}).}
\end{table}
At 2 threads the traffic model's LOVO edge falls far short of the registered $\Delta$AIC margin; at 4 threads the parameter model wins outright, on both criteria. The unscaled raw-byte predictor is far worse (held-out relative error 4.34 at cap 2, 4.97 at cap 4), so the contest is between the fitted models ---
and even fitted, the traffic model fails to justify its extra
structure.
With only six stage-1-valid variants, these point estimates carry wide uncertainty, quantified by a cluster bootstrap over variants (2{,}000 resamples): cap 2
$\Delta$AIC median $-0.91$, 95\% interval $[-28.8, +37.3]$,
$P(\Delta\text{AIC}>10)=0.16$; cap 4 median $-8.90$, interval
$[-138.4, +7.29]$, $P(>10)=0.02$. Read the verdict for what
it is: a verdict under the registered criteria, no significance claim attached. At
cap 2, nothing shows the traffic model better by anything
approaching the registered margin; at cap 4, point estimates and
bootstrap both favour parameters.
A post-hoc sensitivity (requested in the 2026-11-04 panel review; not part of the
frozen criteria) excludes the saturated 16384 context, refitting the identical contest on $C\in\{4096,6144,8192\}$ (98/95 reps).
The cap-2 outcome is unchanged ($\Delta$AIC $+3.70$, still below the
$>10$ margin), but the cap-4 outcome reverses: traffic LOVO 0.207
vs params 0.273, $\Delta$AIC $+20.27$ --- the traffic model wins
there. None of this touches the frozen verdict; the sensitivity localises the failure to the longest context, where saturation and contention peak. Whether it is a true context-scaling
effect or contention noise at 16384, the banked data cannot say.
The amended 10-variant contest, the FLOP-count comparator arm
(Amendment 2 item B1), and the equivalence/power analyses
(Amendments 1--2, A4) consume the amended timing (campaign K2,
cancelled incomplete; the protocol note in Section 3); under A3
their primary set is \{V-depth-lite\}, so they would be reported
as fits including all variants, not as the pre-registered
primary contest. They remain unrun, and no substitute data
are used.
\begin{table}[ht]
\centering\footnotesize\setlength{\tabcolsep}{4pt}
\begin{tabular}{@{}lllllll@{}}
\toprule
Model & Cap & $\beta_0$ & $\beta_1$ & $\beta_2$ & $\beta_3$ & $R^2$ \\
\midrule
$M_P$ & 2 & $-$6.638 & 2.761 & 0.890 & $-$0.347 & 0.412 \\
$M_B$ & 2 & 12.884 & $-$8.433 & 0.519 & --- & 0.413 \\
$M_P$ & 4 & 3.273 & $-$1.287 & $-$0.365 & 0.167 & 0.608 \\
$M_B$ & 4 & 5.790 & $-$3.628 & 0.383 & --- & 0.542 \\
\bottomrule
\end{tabular}
\caption{Fitted models ($y=\log$ CPU-s/token). $M_P$: $y\sim1+P+S+
P\cdot S+L$ ($\beta_1=P$ in $10^9$, $\beta_2=S=(ctx+64)/10^3$,
$\beta_3=P\cdot S$); $M_B$: $y\sim1+B_W+B_{KV}(S)+L$ ($\beta_1=B_W$
in $10^9$ bytes, $\beta_2=B_{KV}$). Layer coefficient $\approx$0
in all four fits ($M_P$ cap 2: 0.001; $M_B$ cap 2: 0.008; $M_P$ cap 4:
$-$0.003; $M_B$ cap 4: $-$0.003). Coefficients are not
interpreted: signs flip across caps (collinearity in a
six-variant fit). Omitted from the traffic model: KV writes,
activation traffic, dequantisation cost, kernel efficiency.
AIC: cap 2 $-328.58$ (P) vs $-330.78$ (B); cap 4 $-507.04$ (P)
vs $-487.94$ (B).}
\end{table}

\subsection{Absolute times}
Absolute times run 0.22--1.38 CPU-s/token. V-depth-lite, the
only D4-valid variant, is slowest at most cap-2 cells (1.376
CPU-s/token at 16384, $\pm$1.164), so validity and speed do not
coincide; the fastest cells belong to V-mixed and V-depth-alt.
Three cells are variance-inflated (caption); the paired ratios
of Section~\ref{sec:hurdle} govern the two-candidate comparison.
\begin{table}[ht]
\centering\scriptsize\setlength{\tabcolsep}{3pt}
\begin{tabular}{@{}lllllllll@{}}
\toprule
Variant & \multicolumn{4}{c}{Cap 2 (CPU-s/token)} & \multicolumn{4}{c}{Cap 4 (CPU-s/token)} \\
 & 4096 & 6144 & 8192 & 16384 & 4096 & 6144 & 8192 & 16384 \\
\midrule
V-width & 0.257 & 0.299 & 0.558 & 0.818 & 0.250 & 0.310 & 0.373 & 0.766 \\
V-depth & 0.254 & 0.474 & 0.415 & 1.187 & 0.244 & 0.500 & 0.372 & 0.983 \\
V-depth-lite & 0.547 & 0.520 & 0.457 & 1.376 & 0.269 & 0.332 & 0.414 & 0.671 \\
V-depth-alt & 0.261 & 0.346 & 0.397 & 0.619 & 0.244 & 0.309 & 0.379 & 0.592 \\
V-KV & 0.262 & 0.323 & 0.380 & 0.918 & 0.245 & 0.303 & 0.355 & 0.669 \\
V-mixed & 0.230 & 0.280 & 0.415 & 0.655 & 0.219 & 0.255 & 0.310 & 0.625 \\
\bottomrule
\end{tabular}
\caption{Cell means, CPU seconds per token ($n=5$ reps per cell),
frozen local timing. 95\% CIs are in the supplementary records;
three cells are variance-inflated (V-width cap-2 8192,
$\pm$0.338; V-depth cap-2 6144, $\pm$0.164; V-depth-lite cap-2
16384, $\pm$1.164); the paired ratios of Section~\ref{sec:hurdle} govern the two-candidate comparison; contexts below 4096 are in the supplementary records; the amended campaign's partial cells are likewise in the supplementary records (campaign K2, incomplete).}
\end{table}

\subsection{Selection regret (descriptive)}
\label{sec:regret}
Regret against the measured fastest primary variant, per cell
(frozen cell means): V-mixed's regret is 0\% at 4096--6144 (both
caps) and at cap-4 8192, reaching 9.4\% at cap-2 8192 (fastest
V-KV) and 5.5--5.8\% at 16384 (fastest V-depth-alt). V-width's
regret is 6.6--47.0\%. Mean regret over $C\geq4096$ is 3.8\%
(V-mixed) vs 24.4\% (V-width) at cap 2 and 1.4\% vs 21.3\% at
cap 4. V-depth-alt, fastest at 16384, is one of the six
primary variants, not an amendment variant. Regret rewards picks rather than models, and V-mixed also
removes five layers, so its speed cannot be credited to its
traffic advantage alone. The registered verdict rests on
Sections~\ref{sec:hurdle}--\ref{sec:contest}; regret shows the
practical stakes were asymmetric: following the traffic model
would have done no harm \emph{here} --- but a model that predicts
neither held-out allocations nor its own margin gives no reason
to expect a repeat elsewhere.

\subsection{Model quality: frozen scorer void; amended re-score reported}
\label{sec:quality}
The A2 quality evaluation ran on Kaggle for all 12 models (40
committed tasks: extraction, arithmetic, format, summarise; greedy
decoding); its frozen scores are void and are cited nowhere as
evidence. The defect is in the instrument, not in any model: tasks
were served by raw completion with no chat template and no stop
strings, so the instruct model answers and then continues
generating, while the frozen scorer assumes the output ends at
the answer --- ``exact'' compares the whole output, ``number''
takes the last integer, and the JSON-array match is greedy across
the continuation. On BASE the first token is correct on 12/12
extraction tasks and the first number on 8/10 arithmetic tasks,
yet the frozen scorer records 0/12 and 0/10 (its 14 of 40 passes
all come from format and summarise); a local re-run reproduces
Kaggle's frozen 35.0 exactly, so the defect is in the design, not
the environment. Amendment 4 (2026-11-05) froze a first-answer
re-scoring rule (first line for exact answers, first integer,
first paragraph for word-count/keywords, first bullet block,
non-greedy JSON match), now executed on the retrieved frozen K1
outputs (the raw pull landed; re-scoring the pulled outputs with
the frozen scorer reproduces every stored score exactly).
\begin{table}[ht]
\centering\scriptsize\setlength{\tabcolsep}{3pt}
\begin{tabular}{@{}lrrrl@{}}
\toprule
Model & Amend. & Frozen & Pass & Ext/Ari/Fmt/Sum (\%) \\
\midrule
BASE & 85.00 & 35.00 & 34/40 & 91.7/80.0/70.0/100.0 \\
V-depth-lite & 55.00 & 25.00 & 22/40 & 91.7/10.0/20.0/100.0 \\
V-shallow & 35.00 & 32.50 & 14/40 & 75.0/0.0/20.0/37.5 \\
V-depth & 30.00 & 7.50 & 12/40 & 75.0/0.0/0.0/37.5 \\
V-KV-lite & 27.50 & 12.50 & 11/40 & 33.3/10.0/30.0/37.5 \\
V-front & 7.50 & 5.00 & 3/40 & 0.0/0.0/0.0/37.5 \\
V-KV & 2.50 & 5.00 & 1/40 & 0.0/0.0/0.0/12.5 \\
V-mid & 2.50 & 2.50 & 1/40 & 0.0/0.0/0.0/12.5 \\
V-depth-alt & 0.00 & 0.00 & 0/40 & 0.0/0.0/0.0/0.0 \\
V-kv2 & 0.00 & 0.00 & 0/40 & 0.0/0.0/0.0/0.0 \\
V-mixed & 0.00 & 0.00 & 0/40 & 0.0/0.0/0.0/0.0 \\
V-width & 0.00 & 0.00 & 0/40 & 0.0/0.0/0.0/0.0 \\
\bottomrule
\end{tabular}
\caption{Amended first-answer re-score (Amendment 4) of the
frozen K1 outputs, mean points of 100, beside the void frozen
scores. Quality is a context axis only, never a verdict
criterion.}
\end{table}
BASE scores 85.00 (34/40); the sole D4-valid variant,
V-depth-lite, is also the highest-scoring pruned variant (55.00);
four variants score 0.00 even under first-answer scoring. The
earlier ``37/40'' figure was a hand diagnostic on locally
regenerated outputs under looser readings, superseded as a
number by this re-score (the verbatim rule gives 34/40 on the K1
outputs; the reconciliation is itemised in the supplementary
record). The corrected rule still detects genuine errors ---
for BASE, ari02, ari06 and fmt09, plus ext05 and fmt05 under the
rule as frozen --- so it is not vacuous. No speed--quality
frontier is claimed on one 40-task suite.

\subsection{Diagnostics and saturation}
\label{sec:diagnostics}
Fresh diagnostic remeasurement of the earliest pool cells did not
reproduce their original slowness: those cells were measured in
the run's highest-contention window (host load 8--10), a load-era
effect noted as a caveat. At 16384, V-depth-alt and V-mixed
saturate at cap 4 (wall-clock tokens/s fall vs cap 2); their cap
comparisons there are saturation-regime numbers. The head-to-head
contrast is load-matched by construction (adjacent paired blocks,
\S\ref{sec:hurdle}); no post-hoc covariate adjustment was
performed (none was pre-registered).

\section{Discussion}
One distinction organises everything below: \emph{selection
regret is not model accuracy.} Section~\ref{sec:regret} shows the
byte-greedy pick was, in hindsight, a low-regret pick; this
section is about why that does not rehabilitate the traffic
model. A model earns its keep by predicting. It predicted
neither held-out allocations (Section~\ref{sec:contest}) nor its
own margins, which it overshot (Section~\ref{sec:hurdle}).
The parameter model should not be oversold either: absolute LOVO
errors are large for \emph{both} models (0.28--0.74 across caps),
so the honest summary is a weak baseline and a challenger that
adds no reliable advantage over it. A good pick justified by an
inaccurate model is luck wearing a plausible story. Five candidate
mechanisms follow, offered as explanations to test --- rescues,
they are not.
\textbf{(1) Parameters already price most of the traffic.} At
fixed precision, weight bytes scale as a constant multiple of the
parameter count, and weights dominate traffic except at the
longest contexts. The KV term --- the traffic model's only
independent signal --- varies less across our allocations than
everything else about them. Even at 16384, where KV reaches
$\approx$56\% of modeled base traffic, the predicted advantage
(0.673) missed the observed ratio (0.986) widely.
\textbf{(2) Logical bytes are not measured
traffic.} Cache residency, tensor layout, and operator fusion all change actual memory movement, which a tensor-byte inventory does not see \citep{hape2025}.
\textbf{(3) Kernel efficiency is shape-dependent.} Cut channels
or heads and matrix shapes change, attained efficiency changes
with them --- work volume is the least of it. That is why
measured-latency methods price whole configurations
\citep{ziplm2023} and roofline predictors need learned
efficiency terms \citep{kernelsight2026}.
\textbf{(4) Pruning damages quantisation.} The variants whose
pruned HF models are most damaged (NLL $+$4.6 to $+$5.5) are the
ones whose GGUFs diverge most from their HF twins (KL
2.2--5.2$\times$ BASE's). A byte-pricing model assumes the
exported model is the pruned model; for 10 of 11 allocations
here, it is not. Finally, \textbf{(5) Thread
saturation on a small, oversubscribed pin.} With 4 threads on 2
pinned cores, scheduling and contention move times by more than
the modeled traffic differences in several cells. The evidence
for saturation is narrow (two variants at 16384,
\S\ref{sec:diagnostics}); a contributing mechanism is all we
claim. Practical guidance is narrow: for a one-off pruning
decision, measure candidate latencies directly (as HAPE and
ZipLM do); between the two analytical models, parameter count
is cheaper and, here, no worse.

\section{Limitations}
\label{sec:limitations}
This is a single model family (Llama-3.2-3B), a single runtime
(llama.cpp, Q4\_K\_M, fp16 KV), a single shared 2-core virtual
machine with substantial tenant load variance (per-rep load recorded; no covariate adjustment pre-registered), and two
thread caps on that pin, which may mask traffic-model
benefits a larger multicore setup would show. Only one of eleven
variants passes the amended export-validity check, so the allocation contest cannot be run on a D4-valid primary set; the frozen contest and hurdle are reported with that caveat. Model quality is reported only as a context axis: the frozen quality instrument is void, and its Amendment-4 first-answer re-score (Section~\ref{sec:quality}) supports no speed--quality frontier claim. The amended-timing analyses (including-fit, TOST, power, FLOP arm) remain unrun: campaign K2 was cancelled incomplete.
The outcome is CPU time, not energy (no energy arm was possible in-environment: no RAPL exposure on this VM or on Kaggle), and decode is CPU-bound at fixed thread counts.
The load-era caveat (§\ref{sec:diagnostics}) applies to
the depth allocations' original cells and runs, if anything,
against them, so it gives the parameter model no advantage in the
primary comparison. The pre-registration history is enumerated
here: SPEC v0.3 frozen 2026-10-23; virgin protocol
adopted after a contamination artifact (30 pre-protocol reps
excluded); base snapshot substituted once; V-mixed-shallow dropped before measurement; three
restart-split paired blocks replaced wholesale; Amendments 1--3
(2026-11-04) and Amendment 4 (2026-11-05, quality-scorer
correction). Beyond this regime --- other CPUs, runtimes,
quantisation, batch sizes --- generality is untested; the claim
stops where the regime does.

\section{Conclusion}
The test was pre-registered and the budgets matched; under it,
parameter count remained unbeaten on the frozen contest --- but the
tightened evidence will not let that stand as a clean negative.
Ten of eleven variants fail the export-validity check --- a count robust to relaxed limits (3 of 11 at most, Section~\ref{sec:validity}) --- so the allocation contest has no D4-valid primary set; the quality axis is measured only as context (BASE 85.00, the sole valid variant 55.00, under the corrected scorer). What this study does
establish, twice over, is a boundary: an analytical weight-plus-KV
byte-traffic model improved neither held-out prediction nor
registered selection for CPU decode of pruned Llama-3.2-3B, and
all but one allocation fail the export-fidelity standard under
the joint prune-and-quantise pipeline (applied jointly
throughout, never separated). One qualification is localised: in a post-hoc sensitivity
excluding the saturated longest context (an exclusion requested
in review, not pre-registered), the traffic model wins the
4-thread contest (Section~\ref{sec:contest}). The boundary is a context boundary --- and the pipeline imposes its own: check the export first.

\section*{Statement of AI Use}
The author uses the AI tools Muse (Muse Spark 1.3, Meta) for
design, execution, drafting, and grammatical editing assistance;
Perplexity AI for literature-discovery assistance; and Google
Gemini and Groq-hosted models for independent review assistance.
The author reviews, edits, and attests to the accuracy of all
contents.

\section*{Data availability}
All timing records, the frozen analysis, the paired-block manifest,
the pre-registered specification with its amendment record
(Amendments 1--4), the K1 validity and quality campaign records,
and the analysis scripts are archived with the study and accompany
the submission as supplementary material. Comparators: parameter count and byte
traffic (run); FLOP count (Amendment 2 B1, unrun: campaign K2
was cancelled incomplete); roofline bound and second runtime (omission justified, not run);
named pruning methods (design-space coverage, not re-run);
energy (probed, not possible in-environment).

\bibliographystyle{plainnat}
\begin{thebibliography}{9}\small
\bibitem[Ashkboos et~al.(2024)]{slicegpt2024} Saleh Ashkboos,
Maximilian L. Croci, Marcelo Gennari do Nascimento, Torsten
Hoefler, and James Hensman. SliceGPT: Compress Large Language
Models by Deleting Rows and Columns. ICLR, 2024.
arXiv:2401.15024.
\bibitem[Bassi et~al.(2026)]{bassi2026} Abnash Bassi, Jaylen Wang,
Fiodar Kazhamiaka, Daniel S. Berger, and Akshitha Sriraman.
Store or Recompute? Characterizing the Carbon Tradeoff of KV
Cache Retention in LLM Inference. EMC$^2$ 2026, ASPLOS.
\bibitem[Bi et~al.(2026)]{rooflinebench2026} Zhen Bi, Xueshu Chen,
Luoyang Sun, Yuhang Yao, Qing Shen, Jungang Lou, and Cheng Deng.
RooflineBench: A
Benchmarking Framework for On-Device LLMs via Roofline Analysis.
arXiv:2602.11506, 2026.
\bibitem[Kurtic et~al.(2023)]{ziplm2023} Eldar Kurtic, Elias
Frantar, and Dan Alistarh. ZipLM: Inference-Aware
Structured Pruning of Language Models. NeurIPS, 2023.
\bibitem[Lv et~al.(2025)]{kvpruner2024} Bo Lv, Quan Zhou, Xuanang
Ding, Yan Wang, and Zeming Ma. KVPruner: Structural
Pruning for Faster and Memory-Efficient Large Language Models.
arXiv:2409.11057, 2024; ICASSP 2025.
\bibitem[Ma et~al.(2023)]{llmpruner2023} Xinyin Ma, Gongfan Fang,
and Xinchao Wang. LLM-Pruner: On the Structural Pruning of Large
Language Models. NeurIPS, 2023. arXiv:2305.11627.
\bibitem[Na et~al.(2024)]{iiswc2024} Seonjin Na, Geonhwa Jeong,
Byung Hoon Ahn, Jeffrey Young, Tushar Krishna, and Hyesoon Kim.
Understanding Performance Implications of LLM Inference on CPUs.
IISWC, 2024, pp. 169--180.
\bibitem[Parihar(2026)]{parihar2026} Raj Parihar. IEI: A Composite
Infrastructure Efficiency Index for GPU Inference. EMC$^2$ 2026,
ASPLOS.
\bibitem[Pope et~al.(2023)]{pope2023} Reiner Pope, Sholto Douglas,
Aakanksha Chowdhery, Jacob Devlin, James Bradbury, Anselm
Levskaya, Jonathan Heek, Kefan Xiao, Shivani Agrawal, and Jeff
Dean. Efficiently Scaling Transformer Inference. MLSys, 2023,
pp. 606--624.
\bibitem[Qi et~al.(2026)]{qi2026} Xiangbo Qi, Chaoyi Jiang, and
Murali Annavaram. Fast NF4 Dequantization Kernels for Large
Language Model Inference. EMC$^2$ 2026, ASPLOS.
arXiv:2604.02556.
\bibitem[Salmani \& Soloveychik(2025)]{salmani2025} Mahsa Salmani
and Ilya Soloveychik. LLM Inference Acceleration via Efficient
Operation Fusion. EMC$^2$ 2025, HPCA.
\bibitem[Sun et~al.(2024)]{wanda2024} Mingjie Sun, Zhuang Liu,
Anna Bair, and J. Zico Kolter. A Simple and Effective Pruning
Approach for Large Language Models. ICLR, 2024.
arXiv:2306.11695.
\bibitem[Wang et~al.(2026)]{grasprune2026} Ziyang Wang,
Jiangfeng Xiao, Chuan Xiao, Ruoxiang Li, Rui Mao, and Jianbin
Qin. GRASPrune: Global Gating
for Budgeted Structured Pruning of Large Language Models.
ACL 2026 (arXiv:2604.19398).
\bibitem[Yao et~al.(2026)]{kernelsight2026} Xiteng Yao, Taeho Kim,
Hengzhi Pei, Xinle Liu, Kyle Ulrich, Leonard Lausen, Ashish
Khetan, Xiang Song, George Karypis, and Martin Herbordt.
KernelSight-LM: A
Kernel-Level LLM Inference Simulator. arXiv:2606.28565, 2026.
\bibitem[Zhao et~al.(2025)]{hape2025} Wenqian Zhao, Lancheng Zou,
Zixiao Wang, Xufeng Yao, and Bei Yu. HAPE: Hardware-Aware LLM Pruning
for Efficient On-Device Inference Optimization. ACM Transactions
on Design Automation of Electronic Systems, 30(4), 2025.
doi:10.1145/3744244.
\end{thebibliography}
\end{document}
